\documentclass[tikz,border=8pt]{standalone}
\usepackage[no-math]{fontspec}
\usepackage[brazil]{babel}
\usepackage{amsmath,amssymb}
\IfFileExists{C:/Windows/Fonts/Lato_Regular.ttf}{%
  \setmainfont{Lato}[
    Path = C:/Windows/Fonts/,
    UprightFont = Lato_Regular.ttf,
    BoldFont = Lato_Bold.ttf,
    ItalicFont = Lato_Italic.ttf,
    BoldItalicFont = Lato_BoldItalic.ttf
  ]
  \setsansfont{Lato}[
    Path = C:/Windows/Fonts/,
    UprightFont = Lato_Regular.ttf,
    BoldFont = Lato_Bold.ttf,
    ItalicFont = Lato_Italic.ttf,
    BoldItalicFont = Lato_BoldItalic.ttf
  ]
}{%
  \RequirePackage[default]{lato}
}
\renewcommand{\familydefault}{\sfdefault}

\usepackage{xcolor}
\usepackage{tikz}
\usepackage{pgfplots}
\pgfplotsset{compat=1.18}
\usetikzlibrary{arrows.meta,positioning,shapes.geometric,fit,backgrounds,calc,decorations.pathreplacing,patterns}

% Paleta institucional UEFS / CHF662
\definecolor{uefsNavy}{HTML}{1B2A4A}
\definecolor{uefsTeal}{HTML}{0E7C7B}
\definecolor{uefsForest}{HTML}{2D6A4F}
\definecolor{uefsGold}{HTML}{C4983B}
\definecolor{uefsGoldDark}{HTML}{8F6B1E}
\definecolor{uefsAlert}{HTML}{B02A37}
\definecolor{uefsOrange}{HTML}{D95F02}
\definecolor{uefsInk}{HTML}{22252A}
\definecolor{uefsMuted}{HTML}{5C6470}
\definecolor{uefsWater}{HTML}{2E86AB}
\definecolor{uefsSand}{HTML}{D9B382}
\definecolor{bgNavy}{HTML}{F2F5F9}
\definecolor{bgTeal}{HTML}{E9F5F4}
\definecolor{bgForest}{HTML}{EEF5F0}
\definecolor{bgGold}{HTML}{FCF7EB}
\definecolor{bgAlert}{HTML}{FDF2F2}
\definecolor{bgGray}{HTML}{F8FAFC}
\definecolor{lineGray}{HTML}{CFD8DC}

\tikzset{
  >=Stealth,
  aaCard/.style={draw=#1, fill=white, rounded corners=4pt, line width=1pt, inner sep=6pt, align=center, font=\small},
  aaBox/.style={draw=#1, rounded corners=4pt, line width=1.1pt, inner sep=6pt, align=center, font=\footnotesize\sffamily, text=uefsInk},
  aaTitle/.style={font=\normalsize\bfseries\sffamily, text=uefsNavy, align=center},
  aaSub/.style={font=\footnotesize\sffamily, text=uefsMuted, align=center},
  aaArrow/.style={->, >=Stealth, line width=1.4pt},
  aaDArrow/.style={<->, >=Stealth, line width=1.4pt}
}


\begin{document}
\begin{tikzpicture}[font=\sffamily]

% Título principal bem acima de toda a figura
\node[font=\large\bfseries\sffamily, text=uefsNavy, align=center] at (5.35, 7.15)
  {Estrutura Operacional da Matriz de Leopold (Leopold et al., 1971)};

% =========================================================
% PAINEL ESQUERDO: MATRIZ CRUZADA (Ações x Fatores)
% =========================================================
\begin{scope}[shift={(0,0)}]
  \def\cw{1.36}
  \def\ch{0.64}
  \def\hh{2.30}

  % Rótulo superior das colunas
  \node[font=\small\bfseries\sffamily, text=uefsNavy] at (2.5*\cw, 6*\ch + \hh + 0.35)
    {Ações do Empreendimento (Colunas)};

  % Fundos dos cabeçalhos das 5 colunas (desenhados primeiro)
  \foreach \j in {0,...,4} {
    \fill[bgNavy] (\j*\cw, 6*\ch) rectangle (\j*\cw+\cw, 6*\ch+\hh);
    \draw[lineGray, line width=0.8pt] (\j*\cw, 6*\ch) rectangle (\j*\cw+\cw, 6*\ch+\hh);
  }
  % Rótulos rotacionados depois (sem oclusão e estritamente dentro de cada coluna)
  \foreach \j/\lbl in {
    0/Supressão veg.,
    1/Terraplenagem,
    2/Obras civis,
    3/Drenagem,
    4/Operação%
  } {
    \node[rotate=65, anchor=west, font=\scriptsize\bfseries\sffamily, text=uefsNavy]
      at (\j*\cw+0.44, 6*\ch+0.12) {\lbl};
  }

  % Rótulo vertical à esquerda dos grupos de fatores
  \node[rotate=90, font=\small\bfseries\sffamily, text=uefsNavy] at (-3.25, 3*\ch)
    {Fatores Ambientais (Linhas)};

  % Faixas de grupo (Físico, Biótico, Socioeconômico)
  \fill[uefsNavy]   (-2.85, 4*\ch) rectangle (-2.3, 6*\ch);
  \node[rotate=90, font=\scriptsize\bfseries\sffamily, text=white] at (-2.575, 5*\ch) {Físico};

  \fill[uefsForest] (-2.85, 2*\ch) rectangle (-2.3, 4*\ch);
  \node[rotate=90, font=\scriptsize\bfseries\sffamily, text=white] at (-2.575, 3*\ch) {Biótico};

  \fill[uefsGoldDark] (-2.85, 0*\ch) rectangle (-2.3, 2*\ch);
  \node[rotate=90, font=\scriptsize\bfseries\sffamily, text=white] at (-2.575, 1*\ch) {Socioec.};

  % Linhas de fatores (y de cima para baixo: row 5 = topo, row 0 = base)
  \foreach \r/\lbl/\bg in {
    5/Solo e erosão/bgNavy,
    4/Água superficial/bgNavy,
    3/Flora nativa/bgForest,
    2/Fauna silvestre/bgForest,
    1/Emprego e renda/bgGold,
    0/Conflitos de uso/bgGold%
  } {
    \fill[\bg] (-2.3, \r*\ch) rectangle (0, \r*\ch+\ch);
    \draw[lineGray, line width=0.8pt] (-2.3, \r*\ch) rectangle (0, \r*\ch+\ch);
    \node[anchor=west, font=\scriptsize\bfseries\sffamily, text=uefsInk] at (-2.18, \r*\ch+0.5*\ch) {\lbl};
  }

  % Grade completa 5x6
  \foreach \j in {0,...,4} {
    \foreach \r in {0,...,5} {
      \draw[lineGray, line width=0.8pt, fill=white] (\j*\cw, \r*\ch) rectangle (\j*\cw+\cw, \r*\ch+\ch);
    }
  }

  % Macro para célula diagonalizada: \cellMI{col}{row}{mag}{imp}{color}
  \newcommand{\cellMI}[5]{%
    \fill[#5!12] (#1*\cw, #2*\ch) rectangle (#1*\cw+\cw, #2*\ch+\ch);
    \draw[lineGray, line width=0.8pt] (#1*\cw, #2*\ch) rectangle (#1*\cw+\cw, #2*\ch+\ch);
    \draw[#5, line width=0.9pt] (#1*\cw, #2*\ch) -- (#1*\cw+\cw, #2*\ch+\ch);
    \node[font=\tiny\bfseries, text=#5] at (#1*\cw+0.30*\cw, #2*\ch+0.72*\ch) {#3};
    \node[font=\tiny\bfseries, text=uefsInk] at (#1*\cw+0.72*\cw, #2*\ch+0.28*\ch) {#4};
  }

  \cellMI{0}{5}{-6}{7}{uefsAlert}
  \cellMI{0}{3}{-8}{9}{uefsAlert}
  \cellMI{0}{2}{-6}{7}{uefsAlert}
  \cellMI{1}{4}{-5}{7}{uefsAlert}
  \cellMI{2}{1}{+6}{7}{uefsForest}
  \cellMI{3}{4}{+5}{6}{uefsForest}
  \cellMI{4}{1}{+7}{8}{uefsForest}
  \cellMI{4}{0}{-4}{6}{uefsAlert}

  % Célula destacada (Terraplenagem x Solo e erosão: col=1, row=5)
  \fill[bgAlert] (1*\cw, 5*\ch) rectangle (2*\cw, 6*\ch);
  \draw[uefsAlert, line width=1.8pt] (1*\cw, 5*\ch) rectangle (2*\cw, 6*\ch);
  \draw[uefsAlert, line width=1.1pt] (1*\cw, 5*\ch) -- (2*\cw, 6*\ch);
  \node[font=\tiny\bfseries, text=uefsAlert] at (1*\cw+0.30*\cw, 5*\ch+0.72*\ch) {-7};
  \node[font=\tiny\bfseries, text=uefsInk]   at (1*\cw+0.72*\cw, 5*\ch+0.28*\ch) {8};

  % Borda externa da matriz
  \draw[uefsNavy, line width=1.2pt] (-2.85, 0) rectangle (5*\cw, 6*\ch);
  \draw[uefsNavy, line width=1.2pt] (0, 6*\ch) rectangle (5*\cw, 6*\ch+\hh);

  % Coordenada âncora da célula destacada
  \coordinate (sampleCell) at (2*\cw, 5.5*\ch);
\end{scope}

% =========================================================
% PAINEL DIREITO: ZOOM DA CÉLULA + LEGENDA + AVISO CRÍTICO
% =========================================================
\begin{scope}[shift={(7.9, 0)}]

  % Caixa de Zoom da Célula
  \draw[draw=uefsNavy, fill=bgGray, rounded corners=6pt, line width=1.1pt] (-0.25, 2.85) rectangle (6.65, 5.85);
  \node[font=\footnotesize\bfseries\sffamily, text=uefsNavy] at (3.2, 5.5)
    {Zoom: Célula da Matriz (Terraplenagem $\times$ Solo)};

  % Quadrado ampliado com corte diagonal
  \begin{scope}[shift={(0.2, 3.15)}]
    \fill[bgAlert] (0,0) -- (0, 1.95) -- (1.95, 1.95) -- cycle;
    \fill[bgNavy]  (0,0) -- (1.95, 0) -- (1.95, 1.95) -- cycle;
    \draw[uefsAlert, line width=1.5pt] (0,0) rectangle (1.95, 1.95);
    \draw[uefsNavy, line width=1.4pt] (0,0) -- (1.95, 1.95);
    \node[font=\normalsize\bfseries, text=uefsAlert] at (0.56, 1.42) {-7};
    \node[font=\normalsize\bfseries, text=uefsNavy]  at (1.42, 0.56) {8};
  \end{scope}

  % Explicação dos triângulos
  \node[anchor=west, text width=3.95cm, align=left, font=\scriptsize\sffamily, text=uefsInk] at (2.4, 4.12) {%
    \textbf{\color{uefsAlert}Magnitude (sup. esq.):}\\
    Escala \textbf{$-10$ a $+10$} (intensidade física da alteração)\\[6pt]
    \textbf{\color{uefsNavy}Importância (inf. dir.):}\\
    Escala \textbf{$1$ a $10$} (peso / relevância no contexto)
  };

  % Legenda de Critérios e Sinais
  \node[draw=uefsTeal, fill=bgTeal, rounded corners=6pt, line width=1.1pt,
        anchor=north west, text width=6.45cm, inner sep=8pt, align=left,
        font=\scriptsize\sffamily, text=uefsInk] at (-0.25, 2.55)
    {\textbf{\footnotesize\color{uefsTeal}Critérios e Leitura dos Sinais}\\[4pt]
     $\bullet$ \textbf{\color{uefsAlert}Sinal ($-$):} impacto adverso (negativo)\\[2pt]
     $\bullet$ \textbf{\color{uefsForest}Sinal ($+$):} impacto benéfico (positivo)\\[2pt]
     $\bullet$ Célula em branco: sem interação direta};

  % Caixa de Aviso Crítico
  \node[draw=uefsAlert, fill=bgAlert, rounded corners=6pt, line width=1.3pt,
        anchor=north west, text width=6.45cm, inner sep=8pt, align=left,
        font=\scriptsize\sffamily, text=uefsInk] at (-0.25, 0.35)
    {\textbf{\footnotesize\color{uefsAlert}Limitação Crítica da Matriz:}\\[3pt]
     \textit{Não captura impactos indiretos de 2ª/3ª ordem, cumulativos ou sinérgicos (exige \textbf{Rede de Interação} complementar).}};

\end{scope}

% Seta conectando a célula destacada ao Zoom (passa limpa pelas células em branco)
\draw[->, >=Stealth, line width=1.6pt, uefsAlert, dashed]
  (sampleCell) -- (7.65, 3.52);

\end{tikzpicture}
\end{document}
